% ============================================================
% Kapitel 3: Sprachverarbeitung: Speech-to-Text
% ============================================================

\chapter{Sprachverarbeitung: Speech-to-Text}
\label{sec:sprachverarbeitung}

Aufbauend auf den im vorherigen Kapitel beschriebenen Grundlagen der Transformer-Architektur wird nun die Sprachverarbeitung behandelt. Um die LLM-Technologie möglichst intuitiv in realen Anwendungen einsetzen zu können, war es in dieser Arbeit das Ziel, gesprochene Sprache in Text umzuwandeln (Speech-to-Text, STT). So kann ein Benutzer Befehle oder Fragen direkt per Mikrofon an das System richten, ohne eine Tastatureingabe vornehmen zu müssen. Bevor die eingesetzte STT-Lösung im Detail erläutert wird, wird zunächst eine physikalische Grundlage von Audiosignalen eingeführt, um die nachfolgenden Abschnitte besser einordnen zu können.

\section{Physikalische Eigenschaften von Audiosignalen}

Beim Sprechen versetzt ein Mensch die Luftmoleküle in Schwingung. Durch die Kompressibilität der Luft breiten sich diese als Schallwellen aus, welche physikalisch als orts- und zeitabhängige Schwankung des Luftdrucks $p$ um den atmosphärischen Ruhedruck ($1013,25$\,hPa) beschrieben werden~\cite{Moeser2015}. Die Einheit des Schalldrucks ist Pascal (Pa).
Der Mensch nimmt Schallwellen primär durch die Klangfarbe (hoch/tief) und die Lautstärke wahr. Die physikalischen Äquivalente hierzu sind die Frequenz $f$ für die Klangfarbe und die Amplitude des Schalldrucks für die Lautstärke. Die Frequenz errechnet sich mit der Einheit Hertz (Hz) über die Kehrwert der Periodendauer $T$:
\begin{equation}
    f = \frac{1}{T}
    \label{eq:frequenz}
\end{equation}

Ein Audiosignal, welches nur eine einzige Frequenzkomponente abdeckt (wie der musikalische Referenzton A4 mit 440\,Hz), nennt man \textit{Ton}. Eine Mischung aus Tönen mit ganzzahligen Frequenzverhältnissen (Grund- und Oberschwingungen) wird dagegen als \textit{Klang} bezeichnet~\cite{Mueller2015}. Die menschliche Sprache entsteht demnach aus einer Überlagerung solcher harmonischer Schwingungen, wobei die relevanten Frequenzen typischerweise in einer Bandbreite zwischen 80\,Hz und 12\,kHz angesiedelt sind.

Da empirische Versuche gezeigt haben, dass das menschliche Lautstärkeempfinden logarithmisch funktioniert, ist die lineare Einheit Pascal zur Einordnung ungeeignet. Stattdessen zieht man den logarithmierten Schalldruckpegel $L$ (angegeben in Dezibel, dB) heran.
\begin{equation}
    L = 20 \cdot \log_{10}\!\left(\frac{p}{p_0}\right) \, \text{dB}
    \label{eq:schalldruckpegel}
\end{equation}
Der Faktor bringt den gemessenen Effektivwert des Schalldrucks $p$ ins Verhältnis zur menschlichen Hörschwelle $p_0$, welche bei $2 \cdot 10^{-5}$\,Pa (entsprechend 0\,dB bei 1\,kHz) liegt. Der Vorfaktor von 20 wurde in der Definition bewusst so gewählt, dass eine Pegeldifferenz von 1\,dB genau der relativen Unterschiedsschwelle des Gehörs entspricht, bei der zwei Schalldrücke gerade noch als unterschiedlich laut ermittelt werden können. Die Skala erstreckt sich dabei von der Hörschwelle bis hin zur Schmerzgrenze bei etwa 200\,Pa, was einem extremen Lautstärkeeindruck von 140\,dB entspricht.
\begin{figure}[H]
    \centering
    \includegraphics[width=0.8\textwidth]{physik_sprache.png}
    \caption{Schematische Darstellung des Datenflusses in einer sprachgesteuerten KI-Pipeline}
    \label{fig:spracherkennung}
\end{figure}

Das Eingangssignal ist eine Audiodatei (typischerweise 16\,kHz Mono). Um es
für eine Transformer-Architektur nutzbar zu machen, wird es in ein
\textbf{Log-Mel-Spektrogramm} umgewandelt, eine Zeit-Frequenz-Darstellung.

% ============================================================
\section{Das eingesetzte System}
\label{sec:system}
% ============================================================

Auf Basis der in den vorherigen Abschnitten beschriebenen Theorie wird
in dieser Arbeit eine vollständige Sprachpipeline eingesetzt\cite{glados_personality}, die
ausschließlich lokal ohne Cloud-Anbindung betrieben wird.
Die zwei Kernkomponenten sind:

\textbf{Spracherkennung: Parakeet TDT 0.6B.}\label{sec:parakeet}
\textbf{Parakeet TDT 0.6B} ist ein von NVIDIA entwickeltes
Open-Source-Modell zur automatischen Spracherkennung (ASR) mit ca.\,600 Millionen Parametern.
Es basiert auf der \textbf{Token-and-Duration Transducer} (TDT)-Architektur,
einer Weiterentwicklung des klassischen RNN-Transducers, die sowohl Token-
als auch Dauerschätzungen für jedes Zeitsegment vorhersagt. Die
TDT-Architektur erreicht dadurch eine höhere Erkennungsgenauigkeit bei
gleichzeitig schnellerer Inferenz, da überflüssige Blank-Token effizienter
übersprungen werden können.
Das Modell wurde auf über 65.000 Stunden englischsprachiger Audiodaten trainiert
und wird in dieser Arbeit im ONNX-Format (Open Neural Network Exchange) lokal
auf der IZPL-Hardware ausgeführt. Die ONNX-Runtime ermöglicht eine
hardwareunabhängige Inferenz sowohl auf GPU als auch auf CPU.
Ein gesprochener Satz wie \glqq{}Bewege den Roboter zur Startposition\grqq{}
wird innerhalb weniger hundert Millisekunden transkribiert und als Zeichenkette
an das Sprachmodell übergeben.

\textbf{Sprachmodell: Qwen 3.5 9B.}\label{sec:qwen}
\textbf{Qwen 3.5 9B} ist ein von Alibaba Cloud entwickeltes
Open-Source-Sprachmodell mit ca.\,9 Milliarden Parametern. Es verwendet
die in Abschnitt~\ref{sec:transformer_block} beschriebene
Transformer-Architektur und zeichnet sich durch eine starke
Leistungsfähigkeit im Verhältnis zu seiner Modellgröße aus.

\begin{table}[H]
    \centering
    \small
    \begin{tabular}{l c c c c r}
        \toprule
        Modell & $N$ & $d$ & $h$ & $d_{\text{ff}}$ & Parameter \\
        \midrule
        GPT-2 (small)       & 12 & 768    & 12 & 3.072   & 124\,M   \\
        LLaMA\,7B            & 32 & 4.096  & 32 & 11.008  & 6{,}7\,B \\
        \textbf{Qwen 3.5 9B} & \textbf{36} & \textbf{3.584} & \textbf{28} & \textbf{18.944} & \textbf{$\approx$9\,B} \\
        GPT-3               & 96 & 12.288 & 96 & 49.152  & 175\,B   \\
        \bottomrule
    \end{tabular}
    \caption{Architekturparameter -- Qwen 3.5 9B (fett) ist das in dieser Arbeit
        eingesetzte Modell.}
    \label{tab:modelle}
\end{table}

Im Vergleich zum Standard-Transformer setzt Qwen 3.5 folgende
Anpassungen ein:
\begin{itemize}
    \item \textbf{RMSNorm} (statt LayerNorm) vor jeder Operation
          (vgl. Gleichung~\ref{eq:rmsnorm})
    \item \textbf{SwiGLU}-Aktivierung im FFN (statt ReLU):
          \begin{equation}
              \text{FFN}_{\text{SwiGLU}}(x) =
              \bigl(\text{Swish}(xW_1) \odot (xW_3)\bigr)\,W_2,
              \quad \text{Swish}(z) = z \cdot \sigma(z)
              \label{eq:swiglu}
          \end{equation}
    \item \textbf{RoPE} Positionskodierung (vgl. Gleichung~\ref{eq:rope})
    \item \textbf{Grouped Query Attention} (GQA) für effizientere Inferenz
    \item Maximale Kontextlänge: $T_{\max} = 32.768$ Tokens
\end{itemize}

Das Modell wird über \texttt{ollama}
lokal betrieben, was einen VRAM-Bedarf von ca.\,6\,GB ergibt und
den Betrieb auf der verfügbaren Hardware ohne relevante
Qualitätseinbußen ermöglicht.

% ============================================================
\section{Zusammenfassung}
% ============================================================

Dieses Kapitel hat die Sprachverarbeitung im Kontext der eingesetzten Pipeline vorgestellt.
Zunächst wurden die physikalischen Grundlagen von Audiosignalen erläutert,
gefolgt von der Beschreibung des konkreten Speech-to-Text-Systems.
Anschließend wurde das eingesetzte Sprachmodell beschrieben:

\begin{table}[H]
    \centering
    \small
    \begin{tabular}{l l l}
        \toprule
        Komponente & Eingesetztes Modell & Besonderheit \\
        \midrule
        STT & Parakeet TDT 0.6B (ONNX) & lokal, TDT-Architektur, 600M Parameter \\
        LLM & Qwen 3.5 9B (ollama) & $N{=}36$, $d{=}3584$, $\approx$6\,GB VRAM \\
        \bottomrule
    \end{tabular}
    \caption{Übersicht der eingesetzten Komponenten.}
    \label{tab:zusammenfassung}
\end{table}

Mit dem Wissen über die eingesetzte Sprachpipeline werden im folgenden Kapitel die Hardwareanforderungen für den lokalen Betrieb eines LLMs und der konkrete Laboraufbau am IZPL beschrieben.
